\documentclass[12pt,a4paper]{article}

% --- Packages ---
\usepackage[utf8]{inputenc}
\usepackage[T1]{fontenc}
\usepackage{mathptmx}           % Times font (close to journal style)
\usepackage[english]{babel}
\usepackage{amsmath,amssymb}
\usepackage{graphicx}
\usepackage{booktabs}
\usepackage{multirow}
\usepackage{array}
\usepackage{tabularx}
\usepackage{float}
\usepackage{hyperref}
\usepackage{natbib}
\usepackage[margin=2.5cm]{geometry}
\usepackage{setspace}
\onehalfspacing
\usepackage{xcolor}

% --- Title ---
\title{Formant Reference of the Orchestra:\\
A Quantitative Study of Instrumental Spectral Convergences\\
and the Acoustic Foundations of Orchestral Doublings}

\author{Yan Maresz\\[6pt]
\small Composer, Professor of New Technologies Applied to Composition\\
\small Conservatoire National Sup\'erieur de Musique et de Danse de Paris\\
\small \texttt{y.maresz@cnsmdp.fr}
}

\date{}

\begin{document}

\maketitle

% ============================================================
% ABSTRACT
% ============================================================
\begin{abstract}
This study presents the first systematic quantitative analysis of the spectral
formants of 30 orchestral instruments, based on 5,914 sustained-tone samples
drawn from the SOL2020 database (IRCAM/Orchidea) and a supplementary
collection (Yan\_Adds). We introduce the Principal Formant (\textit{Fp}), a
band-limited energy-weighted spectral centroid that captures an instrument's
timbral identity with stability gains of up to 10.4$\times$ over raw peak
measurements. Using \textit{Fp} and formant-peak analysis, we demonstrate that
classical orchestral doublings can be explained by measurable formant
convergences. The central finding is a convergence cluster in the vocalic zone
/o/--/\aa/ (377--506~Hz), grouping 11~instruments from 4~different families
with inter-instrument distances as small as $\Delta = 0$~Hz (flute--oboe) and
$\Delta = 11$~Hz (bassoon--violin, horn--viola). Nine principles of acoustic
orchestration are formulated and validated against a multi-source concordance
rate of 93\% (27/29 comparable instruments). By connecting the spectral-centroid
literature on timbral blend \citep{Sandell1995,Tardieu2012,Goodchild2018} with
the orchestration-treatise tradition, this study provides a quantitative
framework for both the analysis and the practice of orchestration.

\medskip
\noindent\textbf{Keywords:} instrumental formants, orchestration, spectral analysis,
doublings, spectral centroid, formant convergence, timbral blend, SOL2020, Orchidea
\end{abstract}

% ============================================================
% 1. INTRODUCTION
% ============================================================
\section{Introduction}

Orchestral doublings---the association of two or more instruments playing the
same melodic line or complementary lines---have constituted one of the
fundamental tools of orchestral writing since the eighteenth century. The
treatise literature, from Berlioz \citep{Berlioz1844} through Rimsky-Korsakov
\citep{RimskyKorsakov1908}, Koechlin \citep{Koechlin1943}, Piston
\citep{Piston1955}, and Adler \citep{Adler1982}, prescribes these associations
on the basis of accumulated auditory experience, without systematic acoustic
quantification.

In parallel, the psychoacoustic study of timbral blend has established that the
spectral centroid plays a central role in determining how well instrument pairs
fuse perceptually. In his seminal work, Sandell demonstrated that blend worsens
as a function of composite centroid height and of centroid difference between
instruments, accounting for up to 63\% of the variance in blend ratings for
dyads separated by a minor third \citep{Sandell1991,Sandell1995}. Subsequent
research has confirmed and extended these findings: Kendall and Carterette
showed that blend can be predicted from proximity in perceptual timbre space
\citep{Kendall1993}; Tardieu and McAdams studied blend perception for dyads of
impulsive and sustained sounds \citep{Tardieu2012}; and Lembke and McAdams
demonstrated that local spectral-envelope features, including formant peaks and
their frequency relationships, predict up to 85\% of the variance in timbral
blend \citep{Lembke2017}. More recently, Goodchild and McAdams proposed a
comprehensive taxonomy of orchestral grouping effects grounded in auditory scene
analysis principles \citep{Goodchild2018,McAdams2022}.

Despite this body of work, a gap remains between the psychoacoustic literature
on timbral blend---which typically studies a small number of instrument pairs
under controlled conditions---and the orchestration-treatise tradition, which
covers the full orchestra but lacks quantitative grounding. The present study
aims to bridge this gap by providing a systematic spectral-formant analysis
of 30 orchestral instruments, identifying the acoustic convergences that
underlie the doublings prescribed by the treatise literature.

The contributions of this paper are threefold: (1)~a comprehensive formant
database covering 30~instruments across 5,914~samples, validated against four
independent academic sources; (2)~the introduction of the Principal Formant
(\textit{Fp}), a band-limited spectral centroid that provides dramatically more
stable timbral characterization than raw peak extraction for instruments with
wide tessitura; and (3)~the identification of a cross-family convergence cluster
in the /o/--/\aa/ vocalic zone that quantitatively explains the most valued
doublings in the orchestral literature.

% ============================================================
% 2. RELATED WORK
% ============================================================
\section{Related Work}

\subsection{Timbre descriptors and the spectral centroid}

The spectral centroid---the center of gravity of the power spectrum---has been
identified as one of the most perceptually relevant timbre descriptors since the
pioneering multidimensional-scaling studies of Grey \citep{Grey1977}. McAdams
et al.\ confirmed its correlation with the ``brightness'' dimension of timbre
spaces \citep{McAdams1995}. The Timbre Toolbox \citep{Peeters2011}, developed
at IRCAM, provides a standardized framework for computing spectral centroid and
related descriptors from musical signals. In the MPEG-7 standard, the spectral
centroid is included as a core audio descriptor \citep{MPEG7}.

A key distinction should be drawn between the \emph{global} spectral centroid
(computed over the full spectrum) and \emph{band-limited} centroid measures. The
global spectral centroid captures overall brightness but does not isolate the
formant structure of an instrument's resonant body. For instruments with wide
tessitura, where the fundamental frequency sweeps across a large range, the
global centroid tracks pitch rather than timbre. The \textit{Fp} measure
introduced in this study addresses this limitation by computing the centroid
within an instrument-specific frequency band optimized to capture the principal
resonance (see Section~\ref{sec:fp}).

\subsection{Formants in musical acoustics}

The concept of formants---fixed resonance regions of the instrument body that
shape the spectral envelope independently of the played pitch---originates in
speech acoustics \citep{Fant1960} and was extended to musical instruments by
Backus \citep{Backus1969}, Meyer \citep{Meyer2009}, and Gieseler et al.\
\citep{Gieseler1985}. Meyer's extensive measurements established vocalic
correspondences for instrumental timbres, describing the horn's timbre as
``/o/-like'' and the oboe's as ``/a/-like.'' However, these characterizations
remained qualitative, based on listening impressions rather than systematic
spectral analysis across large corpora.

Fletcher and Rossing provide the physical acoustics foundation, showing how the
resonant properties of instrument bodies (air columns, string-body coupling,
bell flare) create fixed spectral-envelope features \citep{Fletcher1998}. For
bowed strings, the \emph{Hauptformant} (principal formant) described by Meyer
corresponds to the main body resonance and is relatively stable across the
instrument's range---a property we quantify systematically in this study.

\subsection{Timbral blend in orchestration}

Sandell's doctoral thesis and subsequent publication established the empirical
foundation for understanding timbral blend \citep{Sandell1991,Sandell1995}.
His three experiments with naturally sounding instrument dyads showed that: (a)
instruments with lower spectral centroids blend better; (b) smaller centroid
differences between instruments improve blend; and (c) similarity in attack
characteristics and amplitude-envelope correlation contribute to blend. These
factors accounted for 51\% (unisons) to 63\% (minor thirds) of the variance in
blend ratings.

More recent work by Lembke and McAdams extended this approach to spectral-envelope
descriptions, demonstrating that formant-level features---including the frequency
relationships between spectral maxima, their prominence, and auditory-system
constraints---predict approximately 85\% of blend variance
\citep{Lembke2015,Lembke2017}. This finding motivates the present study's focus
on formant convergence as a predictor of orchestral blend.

The ACTOR project (Analysis, Creation, and Teaching of ORchestration) at McGill
University has produced a taxonomy of orchestral grouping effects grounded in
auditory scene analysis \citep{Goodchild2018,McAdams2022}, distinguishing between
timbral augmentation (a dominant timbre embellished by others) and timbral
emergence (a new timbre not identifiable as any constituent). Our formant-convergence
framework provides acoustic correlates for these perceptual categories: convergent
formants ($\Delta \leq 80$~Hz) predict augmentation/fusion, while divergent
formants ($\Delta > 200$~Hz) predict complementarity or emergence.

\subsection{Computational orchestration}

The Orchidea system \citep{Cella2020}, developed from the SOL database at IRCAM, uses
spectral matching to find combinations of instruments that approximate a target
sound. Our work complements this approach by providing a formant-level analysis
that explains \emph{why} certain combinations work perceptually, rather than
finding combinations through brute-force spectral matching.

% ============================================================
% 3. CORPUS AND SOURCES
% ============================================================
\section{Corpus and Sources}
\label{sec:corpus}

\subsection{Primary data}

The corpus comprises two complementary sources:

\begin{itemize}
\item \textbf{SOL2020 (IRCAM/Orchidea)}---12 solo instruments, 2,391 sustained-tone
samples. The SOL2020 database provides pre-computed spectral envelopes
(\texttt{specenv}, 1,024 bins) derived from recordings made in IRCAM's
anechoic chamber at 44,100~Hz / 24-bit. Instruments covered: piccolo, flute,
oboe, clarinet in B$\flat$, bassoon, horn in F, trumpet in C, trombone, tuba,
alto saxophone, violin, viola, cello, and double bass.

\item \textbf{Yan\_Adds}---18 additional instruments (string ensembles, auxiliary
woodwinds, low brass, muted variants), 1,646 samples sourced from the Vienna
Symphonic Library (VSL), Con Timbre, and personal recordings. All samples were
recorded in controlled studio conditions at 44,100~Hz or 48,000~Hz (resampled
to 44,100~Hz) and normalized to $-1$~dBFS peak. This collection extends the
corpus to instruments and configurations not covered by SOL2020, including
bass flute, contrabass flute, English horn, E$\flat$ clarinet, bass clarinet,
contrabass clarinet, bass trombone, bass tuba, contrabass tuba, contrabassoon,
and string sections (violin ensemble, viola ensemble, cello ensemble, double
bass ensemble), as well as muted variants for strings, horn, trumpet,
trombone, and tuba.
\end{itemize}

In total, 487 data lines from SOL2020 and 46 from Yan\_Adds (CSV v3 format)
cover the 30 instruments analyzed. All data and analysis scripts are available
in the accompanying repository.\footnote{\url{https://github.com/zseramnay/Orchestration}}

\subsection{Validation sources}

Four independent academic references serve as cross-validation:

\begin{enumerate}
\item \textbf{Backus} \citep{Backus1969}---targeted measurements of brass and
woodwind formants.
\item \textbf{Gieseler et al.} \citep{Gieseler1985}---detailed measurements with
explicit vocalic associations for 13 instruments.
\item \textbf{Meyer} \citep{Meyer2009}---acoustic analysis with vocalic color
descriptions and defined vocalic zones.
\item \textbf{McCarty/CCRMA} \citep{McCarty2003}---directional reference with
acknowledged accuracy limitations.
\end{enumerate}

The global concordance rate is 93\% (27 out of 29 comparable instruments). The two
discordances concern the piccolo (Backus cites $\sim$3,400~Hz, likely corresponding
to the radiation peak rather than F1) and the trumpet (high F1 variability across
sources due to dynamic-dependent spectral behavior).

% ============================================================
% 4. METHODOLOGY
% ============================================================
\section{Methodology}
\label{sec:methodology}

\subsection{Measurement scope}

All formant values are measured on sustained tones (\emph{sustained}) in
\emph{ordinario} playing, excluding attack transients, extended techniques, and
extreme dynamics. Included techniques: normal, arco, sustain, ordinario,
vibrato, non-vibrato, flautando, from \emph{pp} to \emph{ff}. Excluded
techniques: pizzicato, col legno, tremolo, trills, artificial harmonics,
multiphonics, flutter tongue, mutes (analyzed separately), glissandi,
portamento. These exclusions ensure that the extracted formants reflect the
structural resonances of the instrument rather than transient spectral artifacts.

\subsection{Formant extraction pipeline}

The extraction pipeline uses the following parameters:

\begin{itemize}
\item \textbf{FFT size:} 4,096 samples at 44,100~Hz, yielding a frequency resolution
of $\sim$10.8~Hz/bin.
\item \textbf{Analysis range:} 100--3,000~Hz.
\item \textbf{Input data:} Pre-computed spectral envelopes (\texttt{specenv}, 1,024
bins) from the Orchidea pipeline, which applies Hann windowing and smoothing
upstream.
\item \textbf{Peak detection:} \texttt{scipy.signal.find\_peaks} with the following
parameters: \texttt{height}~$= -30$~dB relative to the regional maximum,
\texttt{distance}~$= 7$ bins ($\approx$70~Hz minimum inter-peak distance),
\texttt{prominence}~$= 3$~dB.
\end{itemize}

A critical methodological choice is the detection of the \emph{most prominent}
peak (highest amplitude) rather than the \emph{first} peak (lowest frequency).
This captures the perceptually dominant formant and avoids low-frequency
secondary resonances that may be acoustically present but perceptually
subordinate.

The pipeline version (v22) was validated by round-trip testing: all 16 instrument
CSV files were regenerated and compared to reference values, achieving
$\Delta = 0$~Hz across all entries.

\subsection{Terminological note}

In this study, $F_1, F_2, \ldots F_6$ designate the principal spectral peaks
extracted by spectral-envelope analysis. These peaks correspond to the main
acoustic resonances of the instrument system and are compared to cardinal vowels
by perceptual analogy, without claiming strict equivalence with vocal-tract
formants in the sense of Fant \citep{Fant1960}. For cylindrical-bore instruments
(clarinet, flute) where no fixed formant exists in the strict sense, $F_1$
designates the dominant spectral-energy zone.

\subsection{The Principal Formant (\textit{Fp})}
\label{sec:fp}

For instruments with wide tessitura, the $F_1$ peak is highly unstable across
registers. The violin, for example, exhibits $F_1$ with $\sigma = 651$~Hz---a
variation exceeding one octave. This instability makes $F_1$ unsuitable as a
timbral signature.

The \textit{Fp} (Principal Formant) measure resolves this problem. It is defined
as the energy-weighted spectral centroid within an instrument-specific frequency
band:

\begin{equation}
F_p = \frac{\sum_{i} f_i \cdot A_i}{\sum_{i} A_i}
\label{eq:fp}
\end{equation}

\noindent where $f_i$ is the frequency of bin $i$, $A_i$ its amplitude (linear
scale), and the summation extends over a band chosen per instrument.

\paragraph{Relationship to global spectral centroid.}
The \textit{Fp} differs from the standard spectral centroid \citep{Peeters2011}
in two respects: (1)~it is computed within a restricted frequency band rather
than over the full spectrum, and (2)~the band boundaries are optimized per
instrument to capture the principal body resonance (\emph{Hauptformant} in
Meyer's terminology \citep{Meyer2009}). The global spectral centroid, applied to
the full spectrum, tracks overall brightness and is heavily influenced by the
fundamental frequency and by high-frequency energy distribution. By contrast,
\textit{Fp} isolates the resonant signature of the instrument body, which is the
acoustically relevant feature for predicting timbral blend in doublings.

\paragraph{Band selection.}
The frequency bands were determined per instrument by identifying the region of
maximum spectral-envelope energy across all registers in \emph{ordinario}
playing. Typical bands are 600--1,400~Hz for strings and horn, and
1,000--2,000~Hz for flute, oboe, and clarinet. The band boundaries were set to
encompass the $-6$~dB points of the main spectral-envelope peak averaged across
all samples for a given instrument.

\paragraph{Stability gain.}
Table~\ref{tab:fp_stability} shows the stability improvement of \textit{Fp}
over $F_2$ for four representative instruments. The gain ranges from
4.7$\times$ (clarinet) to 10.4$\times$ (trumpet). Twenty-two of the
30 instruments in the corpus benefit from a validated \textit{Fp} measurement.

\begin{table}[htbp]
\centering
\caption{Stability comparison: \textit{Fp} vs.\ $F_2$ (ordinario).}
\label{tab:fp_stability}
\begin{tabular}{lrrrrc}
\toprule
Instrument & $F_p$ (Hz) & $\sigma_{F_p}$ & $F_2$ (Hz) & $\sigma_{F_2}$ & Gain \\
\midrule
Violin     & 893  & 92  & 1,518 & 651   & 7.1$\times$ \\
Trumpet    & 1,046 & 98  & 1,324 & 1,018 & 10.4$\times$ \\
Clarinet B$\flat$ & 1,412 & 169 & 1,130 & 795 & 4.7$\times$ \\
Horn       & 738  & 112 & 1,106 & 734   & 6.5$\times$ \\
\bottomrule
\end{tabular}
\end{table}

\paragraph{$\sigma(F_p)$ as spectral selectivity indicator.}
The standard deviation $\sigma(F_p)$ indirectly reflects the bandwidth of the
principal resonance: a low $\sigma$ indicates a narrow, well-defined formant
(high $Q$ factor), while a high $\sigma$ indicates a broad, diffuse formant.
The violin ($\sigma = 92$~Hz) thus possesses a more spectrally focused timbre
than the B$\flat$ clarinet ($\sigma = 169$~Hz), consistent with perceptual
experience.

\subsection{Two complementary representations per register}

For each principal instrument, the reference document provides two
complementary representations computed register by register (registers defined
according to traditional instrumental nomenclatures):

\begin{enumerate}
\item \textbf{Mean formant profile:} extraction of $F_1$--$F_6$ on each
individual sample, then computation of the median per register. This profile
gives the reference positions of each spectral peak.

\item \textbf{Vocalic spectral map:} mean of all spectral envelopes in the
register, displayed on a logarithmic frequency axis with superimposed vocalic
zones (/u/, /o/, /a/, /e/) and normalized cepstral envelope. This map shows the
actual mean spectral shape with its resonance zones, valleys, and full vocalic
context.
\end{enumerate}

The two approaches converge for well-defined formants (reed instruments, brass)
but may diverge for instruments with wide tessitura (strings, flutes, clarinets)
where averaging the envelope smooths out individual peaks. The formant profile
gives reference positions; the spectral map gives the complete acoustic context.

This dual register-by-register analysis resolves the problem of spectral
variability in wide-tessitura instruments: rather than relying on a single
uninformative global $F_1$ value, register-by-register examination reveals
formant convergences that appear only when instruments play in a specific
register (see Section~\ref{sec:principles}, Principle~8).

% ============================================================
% 5. VOWEL–FREQUENCY CORRESPONDENCE
% ============================================================
\section{Vowel--Frequency Correspondence}

The vocalic system provides an intuitive framework for describing instrumental
timbre. Following Meyer \citep{Meyer2009} and Gieseler et al.\
\citep{Gieseler1985}, the correspondences are shown in
Table~\ref{tab:vowels}.

\begin{table}[htbp]
\centering
\caption{Vocalic zones and representative instruments ($F_1$ or $F_p$).}
\label{tab:vowels}
\begin{tabular}{llll}
\toprule
Vowel & Frequency (Hz) & Perception & Representative instruments \\
\midrule
/u/ & 200--400 & Depth     & Db (172), Tuba (226), Cbn (226), Tbn (237) \\
/o/ & 400--600 & Fullness  & Va (377), Hn (388), ASax (398), EH (452), \\
    &          &           & Cl (463), Bn (495), Vn (506) \\
/\aa/ & 600--800 & Warmth  & Cl$^{E\flat}$ (678), Fl (743), Ob (743), Tpt (786) \\
/a/ & 800--1,250 & Power   & Hn$_{Fp}$ (738), Vn$_{Fp}$ (893), Tpt$_{Fp}$ (1,046) \\
/e/ & 1,250--2,600 & Clarity & Cl$_{Fp}$ (1,412), Fl$_{Fp}$ (1,535), Ob$_{Fp}$ (1,485) \\
/i/ & $>$2,600 & Brilliance & Picc (1,992 $F_2$), Tpt+harmon (2,358) \\
\bottomrule
\end{tabular}
\end{table}

% ============================================================
% 6. RESULTS BY INSTRUMENT FAMILY
% ============================================================
\section{Results by Instrument Family}

\subsection{Woodwinds (range: 226--2,638~Hz)}

The woodwind family presents the greatest timbral diversity in the orchestra.
Double reeds (oboe, English horn, bassoon, contrabassoon) share a common
formant zone around 900--1,200~Hz with a characteristic /a/ coloring. Clarinets
(cylindrical bore, odd harmonics) have no fixed formant in the strict sense:
the spectral peak follows the fundamental or the 3rd harmonic. Flutes (open
pipe) present a spectrum that decreases approximately linearly above the
fundamental.

Key findings: the English horn ($F_1 = 452$~Hz) and the bassoon
($F_1 = 495$~Hz) fall within the /o/ zone, alongside the horn (388~Hz) and the
violin (506~Hz). The English horn and the B$\flat$ clarinet converge at
$\Delta = 11$~Hz (452 vs.\ 463~Hz). The bassoon and the violin converge at
$\Delta = 11$~Hz (495 vs.\ 506~Hz).

\subsection{Brass (range: 162--2,358~Hz)}

Brass instruments present well-defined formants and strong spectral stability
across dynamics (except the trumpet). The horn in F ($F_1 = 388$~Hz, /o/ zone)
is the most versatile brass instrument for doublings, with unanimous agreement
across all four academic sources. The trumpet in C presents an extremely
variable $F_1$ ($\sigma = 642$~Hz), but its \textit{Fp} at 1,046~Hz is
10.4$\times$ more stable than $F_2$. In the low range, the trombone (237~Hz),
bass tuba (226~Hz), and contrabass tuba (226~Hz) form a tight /u/ cluster
($\Delta \leq 11$~Hz), joined by the contrabassoon (also 226~Hz).

\subsection{Saxophones}

The alto saxophone ($F_1 = 398$~Hz) falls squarely within the /o/--/\aa/
convergence zone, alongside the horn ($\Delta = 10$~Hz) and the viola
($\Delta = 21$~Hz). It is the only saxophone present in the SOL2020 corpus;
soprano, tenor, and baritone saxophones constitute a priority extension.

\subsection{Strings (range: 172--1,518~Hz)}

The string family presents the greatest spectral variability in the orchestra,
compensated by the \textit{Fp} measure. The violin ($F_1 = 506$~Hz, /o/) and
the bassoon ($F_1 = 495$~Hz) converge at $\Delta = 11$~Hz. The cello joins
this convergence via its $F_2 = 506$~Hz ($\approx F_1$ of the bassoon).

The section effect is clearly quantified: the ensemble does not shift $F_1$
but smooths the upper harmonics ($F_2$ of solo violin: 1,518~Hz $\rightarrow$
ensemble: 1,163~Hz, a $-23$\% reduction). The perceptual result is a more
blended timbre---this is the textural difference between a quartet and an
orchestral section. This smoothing effect is consistent with Sandell's finding
that lower composite centroids improve blend \citep{Sandell1995}.

% ============================================================
% 7. CONVERGENCE CLUSTER
% ============================================================
\section{The /o/--/\aa/ Convergence Cluster}

The central result of this study is the identification of a convergence cluster
in the /o/--/\aa/ vocalic zone (377--506~Hz), grouping 11 instruments from
4 different families within a span of only 129~Hz (Table~\ref{tab:cluster}).

\begin{table}[htbp]
\centering
\caption{The 6 instruments closest to the convergence cluster center.}
\label{tab:cluster}
\begin{tabular}{llrr}
\toprule
Instrument & Family & $F_1$ (Hz) & $\Delta$ / Violin \\
\midrule
Viola        & Strings     & 377 & 129~Hz \\
Horn         & Brass       & 388 & 118~Hz \\
Alto sax     & Saxophones  & 398 & 108~Hz \\
English horn & Woodwinds   & 452 & 54~Hz \\
Clarinet B$\flat$ & Woodwinds & 463 & 43~Hz \\
Bassoon      & Woodwinds   & 495 & 11~Hz \\
Violin       & Strings     & 506 & --- \\
\bottomrule
\end{tabular}
\end{table}

The tightest doublings---bassoon\,+\,violin ($\Delta = 11$~Hz),
horn\,+\,viola ($\Delta = 11$~Hz)---correspond precisely to the most valued
associations in the orchestral literature. In terms of Bark-scale distance
\citep{Zwicker1961}, $\Delta = 11$~Hz at 500~Hz corresponds to approximately
0.1~Bark, well within the critical bandwidth for spectral integration.

\subsection{Verified doublings}

Table~\ref{tab:doublings} presents the principal doublings classified by
increasing $\Delta$.

\begin{table}[htbp]
\centering
\caption{Principal verified doublings (selection).}
\label{tab:doublings}
\begin{tabular}{llrrrl}
\toprule
Instr.~A & Instr.~B & $F_1$ A & $F_1$ B & $\Delta$ & Quality \\
\midrule
Flute & Oboe & 743 & 743 & 0 & Quasi-perfect $\star$ \\
Tuba CB & Contrabassoon & 226 & 226 & 0 & Quasi-perfect $\star$ \\
Horn & Viola & 388 & 377 & 11 & Quasi-perfect \\
Eng.\ horn & Cl B$\flat$ & 452 & 463 & 11 & Quasi-perfect \\
Violin & Bassoon & 506 & 495 & 11 & Quasi-perfect \\
Trombone & Tuba & 237 & 226 & 11 & Quasi-perfect \\
Eng.\ horn & Horn & 452 & 388 & 64 & Excellent \\
Trumpet & Oboe & 786 & 743 & 43 & Excellent \\
Oboe & Bassoon & 743 & 495 & 248 & Complementary \\
Trumpet & Trombone & 786 & 237 & 549 & Complementary \\
\bottomrule
\end{tabular}
\end{table}

Register-by-register analysis further reveals convergences invisible in global
data. For example, the B$\flat$ clarinet in the chalumeau register (D3--D4,
$F_1 = 215$~Hz) converges exactly ($\Delta = 0$~Hz) with the trombone in
its low register; the horn in the medium register ($F_1 = 323$~Hz) converges
with the viola in the medium register---acoustically confirming the emblematic
Brahms doubling.

It is essential to stress that these convergences are valid at the unison, at
the octave, or at any other interval. Formant convergence is a property of
timbre, not of pitch: two instruments whose formants coincide will fuse
regardless of the pitch distance between the notes played. This is why the
violin--bassoon and horn--viola doublings function across registers without
losing their acoustic effectiveness.

% ============================================================
% 8. NINE PRINCIPLES
% ============================================================
\section{Nine Principles of Acoustic Orchestration}
\label{sec:principles}

The data analysis leads to the formulation of nine fundamental principles.

\paragraph{1.\ Formant convergence (fusion).}
When two instruments share a similar $F_1$ ($\Delta \leq 80$~Hz), their
timbres fuse: the ear perceives a homogeneous color rather than two distinct
voices. Measured thresholds: $\Delta = 0$ (formant unison), $\Delta \leq 30$~Hz
(quasi-perfect fusion, $\leq 0.3$~Bark), $\Delta$ 30--80~Hz (effective fusion,
0.3--0.7~Bark). This finding is consistent with Sandell's observation that
smaller centroid differences improve blend \citep{Sandell1995} and with
critical-bandwidth theory \citep{Zwicker1961}.

\paragraph{2.\ Spectral complementarity (enrichment).}
When $\Delta > 200$~Hz ($> 1.5$~Bark), spectra complement each other without
mutual masking, producing broadband enrichment inaccessible to either instrument
alone. Examples: trumpet (786~Hz) + trombone (237~Hz), $\Delta = 549$~Hz.

\paragraph{3.\ Section effect (homogenization).}
The transition from soloist to ensemble does not shift $F_1$ but smooths upper
harmonics: $F_2$ of the violin drops by 23\% in section, $F_3$ by 16\%. The
perceptual result is a more blended timbre. This is consistent with the
prediction that lower composite centroids improve blend \citep{Sandell1995}.

\paragraph{4.\ Mute as timbral transposition.}
For strings and horn, the mute systematically lowers $F_1$ by 6 to 40\%. For
brass, the effect depends on mute type: cup and harmon mutes produce the
opposite effect---the harmon mute on trumpet shifts $F_1$ from the /a/ to the
/i/ zone (+200\%), the most spectacular case in the corpus.

\paragraph{5.\ Cross-family formant groupings.}
Formant groupings cross organological boundaries: the /u/ family (double bass,
tuba, contrabassoon, trombone), the /o/--/\aa/ family (viola, horn, alto sax,
English horn, clarinet, bassoon, violin), and the /a/ family (via \textit{Fp}
values) are more relevant than traditional instrument families for predicting
timbral fusions.

\paragraph{6.\ The /o/--/\aa/ zone as a convergence attractor.}
The concentration of 11 instruments within 129~Hz around 377--506~Hz corresponds
to the geometric center of the human perceptual spectrum (speaking voice,
presence band). This convergence is consistent with the hypothesis that the
orchestral palette has been historically shaped around the spectral region of
maximal perceptual sensitivity, though establishing a causal relationship would
require further historical and perceptual investigation.

\paragraph{7.\ Role of \textit{Fp} in wide-tessitura doublings.}
For instruments with wide tessitura (violin, trumpet, clarinet), \textit{Fp}
advantageously replaces $F_1$ or $F_2$ in doubling analysis. The stability gain
reaches 10.4$\times$ for trumpet.

\paragraph{8.\ Register-specific convergences.}
Register-by-register analysis reveals convergences invisible in global data:
$\Delta = 0$~Hz between B$\flat$ clarinet (chalumeau) and trombone (low);
between horn (medium) and viola (medium); between horn (medium) and trumpet
(medium).

\paragraph{9.\ Hierarchy of fusion.}
The data establish a quantitative hierarchy: $\Delta \leq 30$~Hz (quasi-perfect
fusion, $\leq 0.3$~Bark), $\Delta$ 30--80~Hz (effective fusion, 0.3--0.7~Bark),
$\Delta$ 80--200~Hz (proximity, 0.7--1.5~Bark), $\Delta \geq 200$~Hz
(spectral complementarity, $\geq 1.5$~Bark).

% ============================================================
% 9. DISCUSSION
% ============================================================
\section{Discussion}

\subsection{Validity of \textit{Fp}}

The \textit{Fp} band-limited spectral centroid proves to be the most relevant
measure for characterizing the global timbre of a wide-tessitura instrument. It
dramatically surpasses $F_1$ and $F_2$ in stability while corresponding to the
\emph{Hauptformant} described by Meyer for strings \citep{Meyer2009}.

While the global spectral centroid \citep{Peeters2011} captures brightness,
\textit{Fp} isolates the instrument's resonant identity---the frequency region
where the body (or air column) imposes its timbral signature regardless of pitch.
This distinction parallels the difference between the ``speaking fundamental
frequency'' and the ``formant frequencies'' in speech: both are spectral centroids,
but measured over different frequency ranges and capturing different perceptual
attributes.

\subsection{Connection to timbral blend research}

The formant convergences identified in this study provide direct acoustic
correlates for the blend phenomena studied by Sandell \citep{Sandell1995},
Kendall and Carterette \citep{Kendall1993}, and Lembke and McAdams
\citep{Lembke2017}. Specifically:

\begin{itemize}
\item The $\Delta \leq 80$~Hz threshold for fusion corresponds to a Bark-scale
distance of $\leq 0.7$~Bark at 500~Hz, consistent with critical-bandwidth
constraints on spectral integration \citep{Zwicker1961}.

\item The finding that lower-formant instruments (horn, viola, bassoon in the
/o/ zone) participate in the most effective doublings is consistent with
Sandell's finding that blend improves with lower composite centroids.

\item The section effect (smoothing of upper harmonics without $F_1$ shift)
provides a mechanism for the well-known perceptual observation that
orchestral sections blend better than solo instruments.
\end{itemize}

\subsection{Limitations}

Several limitations should be acknowledged:

\begin{enumerate}
\item \textbf{Corpus coverage:} only one saxophone (alto) is included. Soprano,
tenor, and baritone saxophones constitute a priority extension.

\item \textbf{Heterogeneous sources:} the Yan\_Adds samples come from diverse
recording conditions. While all were recorded in studio or anechoic conditions,
variations in microphone placement and room characteristics may introduce
systematic biases. The 93\% concordance rate with independent academic
sources provides indirect validation.

\item \textbf{No perceptual validation:} the link between formant convergence
and perceived blend is supported by the existing psychoacoustic literature
\citep{Sandell1995,Lembke2017} but has not been tested experimentally with
the specific instrument pairs and $\Delta$ values reported here. A listening
experiment measuring blend ratings as a function of $F_1$ distance for the
doublings in Table~\ref{tab:doublings} would be a valuable next step.

\item \textbf{Sustained tones only:} attack transients, extended techniques, and
extreme dynamics, which radically modify the spectrum, are not covered.
Previous research has shown that attack similarity is a significant predictor
of blend \citep{Sandell1995}, suggesting that a complete model of orchestral
fusion should incorporate temporal as well as spectral features.

\item \textbf{Band selection for \textit{Fp}:} the instrument-specific
frequency bands were determined by spectral-envelope analysis and visual
inspection, introducing a degree of subjectivity. An automated optimization
procedure (e.g., minimizing cross-register variance of the centroid) could
improve reproducibility.

\item \textbf{Statistical testing:} the concordance rate (93\%) is a simple
ratio. Future work should include bootstrapped confidence intervals for
$F_1$ and \textit{Fp} estimates and permutation tests for the significance
of the /o/--/\aa/ cluster relative to a null model of randomly distributed
formants.
\end{enumerate}

\subsection{Implications for composition}

The quantitative identification of cross-family formant groupings offers
composers a new tool: the ability to predict timbral fusions and
complementarities before writing, by choosing instruments according to their
vocalic zone rather than their organological family. The threshold of
$\Delta \leq 80$~Hz for fusion and $\Delta \geq 200$~Hz for spectral
enrichment provides operational criteria that complement, rather than replace,
the composer's ear.

% ============================================================
% 10. CONCLUSION
% ============================================================
\section{Conclusion}

This study establishes that classical orchestral doublings rest on measurable
formant convergences. The /o/--/\aa/ convergence cluster (377--506~Hz), grouping
11 instruments from 4 families, constitutes the acoustic core of the symphonic
orchestra. The introduction of \textit{Fp} (band-limited spectral centroid) as a
timbral metric, combined with register-by-register analysis and multi-source
validation (93\% concordance), provides a quantitative framework connecting the
orchestration-treatise tradition with the psychoacoustic literature on timbral
blend.

The complete reference document, including detailed formant profiles, register-by-register
vocalic spectral maps, and convergence matrices for all 30 instruments, is
available online at \url{https://zseramnay.github.io/Orchestration/}.

\subsection*{Data availability statement}
All data (CSV files), analysis scripts (Python), and the complete reference
document are publicly available at
\url{https://github.com/zseramnay/Orchestration}.

\subsection*{Conflict of interest}
The author declares no conflict of interest.

\subsection*{Acknowledgments}
The author thanks the IRCAM Orchidea team for making the SOL2020 database
available.

% ============================================================
% REFERENCES
% ============================================================
\bibliographystyle{apalike}

\begin{thebibliography}{99}

\bibitem[Adler, 1982]{Adler1982}
Adler, S. (1982).
\newblock \emph{The Study of Orchestration}.
\newblock W.W. Norton \& Company, New York.

\bibitem[Backus, 1969]{Backus1969}
Backus, J. (1969).
\newblock \emph{The Acoustical Foundations of Music}.
\newblock W.W. Norton \& Company, New York.

\bibitem[Berlioz, 1844]{Berlioz1844}
Berlioz, H. (1844/1855).
\newblock \emph{Grand trait\'e d'instrumentation et d'orchestration modernes}.
\newblock Schonenberger, Paris.

\bibitem[Cella et~al., 2020]{Cella2020}
Cella, C.-E., Music{\'e}, E., and Music{\'e}, C. (2020).
\newblock OrchideaSOL: a dataset of extended instrumental techniques for
  computer-aided orchestration.
\newblock In \emph{Proceedings of the International Computer Music Conference}.

\bibitem[Fant, 1960]{Fant1960}
Fant, G. (1960).
\newblock \emph{Acoustic Theory of Speech Production}.
\newblock Mouton, The Hague.

\bibitem[Fletcher and Rossing, 1998]{Fletcher1998}
Fletcher, N.~H. and Rossing, T.~D. (1998).
\newblock \emph{The Physics of Musical Instruments}.
\newblock Springer-Verlag, 2nd edition.

\bibitem[Gieseler et~al., 1985]{Gieseler1985}
Gieseler, W., Lombardi, L., and Weyer, R.-D. (1985).
\newblock \emph{Instrumentation in der Musik des 20.\ Jahrhunderts}.
\newblock Breitkopf \& H{\"a}rtel, Wiesbaden.

\bibitem[Goodchild and McAdams, 2018]{Goodchild2018}
Goodchild, M. and McAdams, S. (2018).
\newblock Perceptual processes in orchestration.
\newblock In Siedenburg, K., Saitis, C., McAdams, S., Popper, A.~N., and
  Fay, R.~R., editors, \emph{Timbre: Acoustics, Perception, and Cognition},
  pages 247--271. Springer.

\bibitem[Grey, 1977]{Grey1977}
Grey, J.~M. (1977).
\newblock Multidimensional perceptual scaling of musical timbres.
\newblock \emph{Journal of the Acoustical Society of America}, 61(5):1270--1277.

\bibitem[{ISO/IEC}, 2002]{MPEG7}
{ISO/IEC} (2002).
\newblock {ISO/IEC 15938-4: Information technology---Multimedia content
  description interface---Part 4: Audio}.

\bibitem[Kendall and Carterette, 1993]{Kendall1993}
Kendall, R.~A. and Carterette, E.~C. (1993).
\newblock Identification and blend of timbres as a basis for orchestration.
\newblock \emph{Contemporary Music Review}, 9(1--2):51--67.

\bibitem[Koechlin, 1943]{Koechlin1943}
Koechlin, C. (1943).
\newblock \emph{Trait\'e de l'orchestration} (4 volumes).
\newblock Max Eschig, Paris.

\bibitem[Lembke and McAdams, 2015]{Lembke2015}
Lembke, S.-A. and McAdams, S. (2015).
\newblock Spectral-envelope characteristics as predictors of blend in
  orchestral instrument pairs.
\newblock In \emph{Proceedings of the International Symposium on Musical
  Acoustics (ISMA)}.

\bibitem[Lembke et~al., 2017]{Lembke2017}
Lembke, S.-A., Parker, K., Narmour, E., and McAdams, S. (2017).
\newblock Predicting blend between orchestral timbres using generalized
  spectral-envelope descriptions.
\newblock \emph{Journal of New Music Research}, 46(4):1--18.

\bibitem[McAdams et~al., 1995]{McAdams1995}
McAdams, S., Winsberg, S., Donnadieu, S., De~Soete, G., and Krimphoff, J.
  (1995).
\newblock Perceptual scaling of synthesized musical timbres: Common dimensions,
  specificities, and latent subject classes.
\newblock \emph{Psychological Research}, 58(3):177--192.

\bibitem[McAdams, 2022]{McAdams2022}
McAdams, S. (2022).
\newblock A taxonomy of orchestral grouping effects derived from principles of
  auditory perception.
\newblock \emph{Music Theory Online}, 28(3).

\bibitem[McCarty, 2003]{McCarty2003}
McCarty, M. (2003).
\newblock Vowel space chart for orchestral instruments.
\newblock CCRMA, Stanford University.

\bibitem[Meyer, 2009]{Meyer2009}
Meyer, J. (2009).
\newblock \emph{Acoustics and the Performance of Music}.
\newblock Springer-Verlag, 5th edition.

\bibitem[Peeters et~al., 2011]{Peeters2011}
Peeters, G., Giordano, B.~L., Susini, P., Misdariis, N., and McAdams, S.
  (2011).
\newblock The {T}imbre {T}oolbox: Extracting audio descriptors from musical
  signals.
\newblock \emph{Journal of the Acoustical Society of America},
  130(5):2902--2916.

\bibitem[Piston, 1955]{Piston1955}
Piston, W. (1955).
\newblock \emph{Orchestration}.
\newblock W.W. Norton \& Company, New York.

\bibitem[Rimsky-Korsakov, 1908]{RimskyKorsakov1908}
Rimsky-Korsakov, N. (1908).
\newblock \emph{Principles of Orchestration} (posthumous).
\newblock Edition russe de musique.

\bibitem[Sandell, 1991]{Sandell1991}
Sandell, G.~J. (1991).
\newblock \emph{Concurrent Timbres in Orchestration: A Perceptual Study of
  Factors Determining ``Blend''}.
\newblock PhD thesis, Northwestern University.

\bibitem[Sandell, 1995]{Sandell1995}
Sandell, G.~J. (1995).
\newblock Roles for spectral centroid and other factors in determining
  ``blended'' instrument pairings in orchestration.
\newblock \emph{Music Perception}, 13(2):209--246.

\bibitem[Tardieu and McAdams, 2012]{Tardieu2012}
Tardieu, D. and McAdams, S. (2012).
\newblock Perception of dyads of impulsive and sustained instrument sounds.
\newblock \emph{Music Perception}, 30(2):117--128.

\bibitem[Zwicker and Fastl, 1999]{Zwicker1961}
Zwicker, E. and Fastl, H. (1999).
\newblock \emph{Psychoacoustics: Facts and Models}.
\newblock Springer-Verlag, 3rd edition.

\end{thebibliography}

\end{document}
